\section{Related Work}
\label{sec:related}

\textbf{Pressure signals.} Linux PSI~\cite{weiner2018psi} exposes per-resource stall time from kernel accounting and underpins user-space managers that throttle or kill work by cgroup. Windows has no user-mode equivalent: performance counters report utilisation and queue lengths, which are proxies for, not measurements of, lost foreground time. WSI recovers a PSI-like quantity from a probe's experience and adds attribution without privileges.

\textbf{Interference control in datacentres.} Heracles~\cite{lo2015heracles} and PARTIES~\cite{chen2019parties} co-locate latency-critical and batch work under feedback controllers over cores, cache ways, memory bandwidth, and power. Caladan~\cite{fried2020caladan} reacts at microsecond scale using kernel-bypass signals. Dirigent~\cite{zhu2016dirigent} predicts progress to steer frequency and cache partitioning. Bubble-Up and Bubble-Flux~\cite{mars2011bubbleup,yang2013bubbleflux} use synthetic pressure probes to measure sensitivity. These systems assume server hardware partitioning (CAT, MBA), privileged control, and an instrumented latency-critical service. SmartSwap targets the opposite setting: a stock desktop, no privileges, no partitioning hardware, and a foreground that cannot be instrumented. It is limited to scheduler hints and CPU-rate caps, so it must decide \emph{which} of these weak tools works and whether it actually helped.

\textbf{Desktop controls.} Windows exposes priority classes, I/O priority hints, Job Object CPU-rate limits~\cite{msjobobjects}, and EcoQoS~\cite{msqos}; scheduling internals are described in~\cite{russinovich2017internals}. Commercial ``process balancers'' apply static or heuristic rules without verification or cost accounting. We include the strongest static rule (\texttt{IDLE} priority plus very-low I/O priority) as a baseline.

\textbf{Measurement rigour.} Systems measurements are easily biased by setup details~\cite{mytkowicz2009wrongdata}; rigorous practice calls for repetition, randomisation, and effect sizes with intervals~\cite{georges2007rigorous,kalibera2013rigorous}. \S\ref{sec:lessons} reports how our own first evaluation failed these standards.

